\subsection{Блочный линейный оракул}

Для $D=D_1\times\dots\times D_n$ полный линейный оракул равен
\[
    \argmin_{v\in D}\inner{g}{v}.
\]
Он распадается по блокам:
\[
    v[k]\in\argmin_{u\in D_k}\inner{g[k]}{u}.
\]
DBCFW использует только выбранный набор $I_t^i$:
\[
    v_t^i[k]\in
    \begin{cases}
        \argmin_{u\in D_k}\inner{\tilde g_t^i[k]}{u}, & k\in I_t^i,\\
        \tilde x_t^i[k], & k\notin I_t^i.
    \end{cases}
\]
Таким образом, невыбранные блоки не меняются в направлении оракула.
После FW-шага они также остаются ближе к текущей смешанной точке:
\[
    x_{t+1}^i[k]=\tilde x_t^i[k],
    \qquad k\notin I_t^i,
\]
если записывать $v_t^i[k]=\tilde x_t^i[k]$.

\subsection{Специальные случаи DBCFW}

DBCFW содержит несколько известных схем как частные случаи.
Если $N=1$, то нет децентрализации и коммуникаций.
Алгоритм превращается в стохастический блочно-координатный FW:
\[
    x_{t+1}=(1-\gamma_t)x_t+\gamma_tv_t,
\]
где $v_t$ отличается от $x_t$ только на выбранных блоках.

Если $B=n$, то каждый агент выбирает все блоки.
Алгоритм превращается в DFW с градиентным отслеживанием.

Если одновременно $N=1$ и $B=n$, то DBCFW становится классическим FW.
Именно поэтому DBCFW удобно рассматривать как общий каркас:
\[
    \text{FW}
    \subset
    \text{BCFW}
    \subset
    \text{DBCFW},
    \qquad
    \text{DFW}
    \subset
    \text{DBCFW}.
\]

\subsection{Интерпретация констант в оценке}

В теореме~\ref{thm:dbcfw-rate} константа
\[
    A_\tau=
    3\bar\rho(C_g+LC_p)+\frac{L}{2}\bar\rho^2
\]
содержит два типа ошибок.
Слагаемое
\[
    3\bar\rho(C_g+LC_p)
\]
отвечает за несовпадение локальных направлений с направлением, которое использовал бы
централизованный алгоритм.
Здесь $C_g$ связан с ошибкой градиентного отслеживания, а $LC_p$ возникает из-за того,
что локальные точки отличаются от средней точки.
Слагаемое
\[
    \frac{L}{2}\bar\rho^2
\]
является обычной гладкостной ценой FW-шага.

Параметр
\[
    \tau=\frac{2n}{B}
\]
увеличивается при малом батче.
Если $B$ мал, шаг должен быть осторожнее, потому что один оракул покрывает меньшую часть
координат.
Если $B=n$, то $\tau=2$, и схема ближе к полному DFW.

\subsection{Итоги раздела}

DBCFW формулируется как общий метод для распределенной проекционно-свободной
оптимизации на декартовых множествах.
Главный алгоритмический объект -- локальный блочный линейный оракул, который вызывается
после смешивания точек и градиентного отслеживания.
Главный теоретический результат -- сохранение порядка $O(1/t)$: леммы о консенсусе и
градиентном отслеживании дают нужное убывание ошибок, а теорема использует эти оценки
в рекурсии спуска.
Главная связь с предыдущими главами состоит в том, что стохастический выбор блоков из SOFW
становится частью распределенного FW-каркаса.
